\subsection{Revisiting modular symmetry in magnetized torus and orbifold compactifications --- arXiv:2005.12642}

\textbf{Authors:} Shota Kikuchi, Tatsuo Kobayashi, Shintaro Takada, Takuya H. Tatsuishi, Hikaru Uchida

\paragraph{主な主張}
磁束 $M$ を持つ $T^2$ 上では $|M|$ 個の zero mode と massive mode の波動関数が modular weight $1/2$ の modular form として $\widetilde{SL}(2,\mathbb{Z})$ の表現をなし、twisted・shifted orbifold $T^2/\mathbb{Z}_N$ では twist・shift の固有値によってそれらの表現がより小さな表現へ分解されることを示した \cite{Kikuchi:2020magnetized}。

\paragraph{新規性}
磁化 $T^2$ の zero mode に加えて massive mode の modular transformation まで統一的に扱い、さらに twisted・shifted orbifold における projection を modular representation の分解として体系化した点にある。

\paragraph{位置付け}
磁化トーラスにおける modular flavor symmetry の top-down な起源と orbifold compactification の関係を整理した基礎文献であり、flux を導入した $T^2$ で波動関数・縮退度・有限モジュラー表現がどのように変化するかを考える際の参照点となる。
